%% macros.tex -- packages, method names and notation for the paper.
%% Loaded by main.tex after the kit's preamble.tex and BEFORE hyperref.
%%
%% Already loaded by cvpr.sty, do not load again: times, xspace, xcolor[dvipsnames], graphicx, amsmath, amssymb,
%% booktabs, natbib[numbers,sort&compress], caption, subcaption, url, etoolbox, enumitem, lineno (review mode),
%% and cleveref (\AtEndPreamble, i.e. after hyperref).  cvpr.sty also defines \eg \ie \cf \etc \vs \wrt \etal.
%% The kit's preamble.tex defines \red{..}, \todo{..} and \TODO{..} (red, in the PDF).

%% ------------------------------------------------------------------ packages
%% Engine guard: the kit loads \usepackage{times} without fontenc.  Under XeTeX/LuaTeX (TU encoding) that silently
%% falls back to Latin Modern body text (checked with tectonic 0.17), so only pdfLaTeX gives the CVPR typography.
\usepackage{iftex}
\ifPDFTeX\else
  \PackageError{macros}{Compile with pdfLaTeX (./build.sh or Overleaf's pdfLaTeX).
    XeTeX/LuaTeX replace the template's Times body font by Latin Modern}{Use pdflatex, e.g. ./build.sh}
\fi
\usepackage{microtype}  % pdfTeX: character protrusion + font expansion. Enabled from day one so that the page
                        % count seen while writing is the final one (the kit suggests it to save space).
\usepackage{bm}         % bold math, incl. Greek (\bm{\Phi})
\usepackage{multirow}
\usepackage{pifont}     % \cmark / \xmark in ablation tables
\usepackage{algorithm}      % float for the algorithm box (sections/method.tex, \label{alg:kalsea})
\usepackage[noend]{algpseudocode}  % algorithmicx syntax (\State, \If, \For, \Comment); noend saves lines
\usepackage{siunitx}    % S columns; numbers follow the text font and weight (\bfseries works in S columns)
\sisetup{mode=match, propagate-math-font=true, reset-math-version=false, reset-text-family=false,
  reset-text-series=false, reset-text-shape=false, text-family-to-math=true, text-series-to-math=true,
  group-digits=false}

%% Double-blind hygiene: do not copy file names or info dictionaries of included PDFs into the output.
\ifdefined\pdfsuppressptexinfo \pdfsuppressptexinfo=-1 \fi

%% Shared method colors (figures/src/palette.tex: cOurs, cKalSea, cSeaCache, cReference, ...), so text,
%% captions and tables can use the same colors as the figures.  Skipped silently if the file does not exist.
\InputIfFileExists{figures/src/palette.tex}{}{}

%% ------------------------------------------------------------------ title
%% Shared by main.tex and supp_main.tex (the supplement repeats the title). Candidates: TITLES.md.
\newcommand{\papertitle}{KalSea: Kalman-Corrected Spectral-Evolution-Aware Caching}

%% ------------------------------------------------------------------ method names
%% Use these macros everywhere (text, tables, captions), never the literal names.
%% Naming (owner, 2026-10-23): the full method is KalSea; the predictor alone is "KalSea w/o PCGS". Keep
%% figures/src/figstyle.py (LABEL) and figures/src/fig_qualitative.py (LABEL) in sync.
\newcommand{\ours}{KalSea\xspace}               % full method: SEA gate + Kalman residual forecasting + gap smoothing
\newcommand{\kalsea}{KalSea w/o \pcgs}          % ablation: SEA gate + Kalman residual forecasting, no gap smoothing
\newcommand{\seacache}{SeaCache\xspace}  % SEA gate + residual reuse (Chung et al., CVPR 2026)
\newcommand{\pcgs}{PCGS\xspace}          % predictor-corrector gap smoothing (the new component)
\newcommand{\sea}{SEA\xspace}            % Spectral-Evolution-Aware (filter / gate) of SeaCache
\newcommand{\nfe}{NFE\xspace}            % number of transformer (block-stack) evaluations per image
\newcommand{\flux}{FLUX.1-dev\xspace}

%% ------------------------------------------------------------------ notation
%% One symbol per quantity, everywhere.  i: sampler step (1..N); a < b: consecutive computed steps; a+1..b-1: gap.
%%   \lat_i  latent x_i            \vel_i  velocity             \sigma_i  noise level (Euler: x_{i+1} = x_i + (\sigma_{i+1}-\sigma_i) v_i)
%%   \hin_i, \hout_i  input / output of the transformer block stack \blocks (image tokens);  \cond_i  conditioning (temb)
%%   \res_i = \hout_i - \hin_i  block-stack residual;  \resf_i  forecast or smoothed residual used on skipped steps
%%   \Emb(.)  input embedding (x_embedder);  \Out(., \cond_i)  output layers (norm_out + proj_out), always run
%%   \modin_i  timestep-modulated input of the first block;  \seaop(.)  SEA filter;  \gdist_i  SEA-filtered
%%             relative L1 change;  \delta  gate threshold (skip while the accumulated \gdist stays below \delta)
%%   Kalman, per element (or per DCT band k):  state \kstate_i = [m_i, \nu_i]^T (value, derivative), warped time s,
%%             \ds_i step, \rho = exp(-\ds/\tau), \kA transition, \kP covariance, \kQ = q \kappa [[ds^3/3, ds^2/2],
%%             [ds^2/2, ds]] process noise (\kappa adapted online), R measurement noise, z_b observation,
%%             e_b = z_b - \prior{m}_b innovation, \kPhi(b,i) transition from i to b, g_i smoother gain
%%   \prior{m}_i  predicted (prior);  \post{m}_i  filtered (posterior);  \smooth{m}_i  RTS-smoothed
\newcommand{\lat}{\mathbf{x}}
\newcommand{\vel}{\mathbf{v}}
\newcommand{\hin}{\mathbf{h}^{\mathrm{in}}}
\newcommand{\hout}{\mathbf{h}^{\mathrm{out}}}
\newcommand{\res}{\mathbf{r}}
\newcommand{\resf}{\tilde{\mathbf{r}}}
\newcommand{\cond}{\mathbf{c}}
\newcommand{\blocks}{\mathcal{F}}
\DeclareMathOperator{\Emb}{Emb}
\DeclareMathOperator{\Out}{Out}
\DeclareMathOperator{\seaop}{SEA}
\DeclareMathOperator{\DCT}{DCT}
\newcommand{\modin}{\mathbf{I}}
\newcommand{\gdist}{d}
\newcommand{\kstate}{\mathbf{s}}
\newcommand{\kA}{\mathbf{A}}
\newcommand{\kP}{\mathbf{P}}
\newcommand{\kQ}{\mathbf{Q}}
\newcommand{\kPhi}{\bm{\Phi}}
\newcommand{\ds}{\Delta s}
\newcommand{\prior}[1]{#1^{-}}
\newcommand{\post}[1]{#1^{+}}
\newcommand{\smooth}[1]{#1^{\mathrm{s}}}
\newcommand{\E}{\mathbb{E}}
\DeclareMathOperator*{\argmin}{arg\,min}
\DeclareMathOperator*{\argmax}{arg\,max}

%% ------------------------------------------------------------------ tables
\newcommand{\best}[1]{\textbf{#1}}          % best value in a column
\newcommand{\second}[1]{\underline{#1}}     % second best
\newcommand{\nodata}{--}                    % intentionally blank cell (TFLOPs / latency until measured)
\newcommand{\tbd}{\textcolor{red}{TBD}\xspace} % result not available yet (video); must be gone before submission
\newcommand{\cmark}{\ding{51}}
\newcommand{\xmark}{\ding{55}}

%% ------------------------------------------------------------------ placeholders for parallel work
%% \maybeinput{tables/foo}  inputs tables/foo.tex if it exists, else draws a red box.
%% \maybegraphics[opts]{figures/foo.pdf}  includes the figure if it exists, else draws a red box of the same width.
%% Missing files are logged as "MISSING-FILE: <name>"; build.sh / tools/pdfcheck.py report them.
\newcommand{\missingbox}[2]{%
  \typeout{MISSING-FILE: #2}%
  \fbox{\parbox[c][1.5cm][c]{\dimexpr#1-2\fboxsep-2\fboxrule\relax}{\centering\footnotesize\color{red}%
    missing: \texttt{\detokenize{#2}}}}}
\newcommand{\maybeinput}[1]{\IfFileExists{#1.tex}{\input{#1.tex}}{\IfFileExists{#1}{\input{#1}}{%
  \missingbox{\linewidth}{#1}}}}
\newcommand{\maybegraphics}[2][width=\linewidth]{\IfFileExists{#2}{\includegraphics[#1]{#2}}{%
  \missingbox{\linewidth}{#2}}}
